% I skabelonen er der tre primære ting at ændre: linjerne øverst, brødteksten og teksten i højre kolonne.

% VIGTIGT! For at dokumentet kan compile, skal man benytte XeLaTeX eller LuaLaTeX som compiler. Dette kan gøres i Overleaf ved at gå til Menu -> Settings -> Compiler -> Vælg XeLaTeX/LuaLaTeX

% Brugerspecifikke variable fra YAML-header
\newcommand{\Titel}{$title$}   % Rapportoverskrift (1-2 linjer)
\newcommand{\Undertitel}{$undertitel$}   % Rapportunderoverskrift (valgfri)
\newcommand{\Dato}{$if(dato)$$dato$$else$\today$endif$}   % Tekst/dato på forsiden
\newcommand{\Afsender}{$if(afsender)$$afsender$$else$Rektoratets Stab · Analyse$endif$}   % Afsenderlinje på forsiden
\newcommand{\Indholdstitel}{$if(indholdstitel)$$indholdstitel$$else$Indhold$endif$}
 
\documentclass[a4paper,12pt]{article}
 
\usepackage{geometry}
\usepackage{xcolor}
\usepackage{microtype}
\usepackage{graphicx}
\usepackage{tikz}
\usepackage{setspace}
\usepackage{fontspec}
\usepackage{hyperref}
\hypersetup{
    colorlinks=true,
    linkcolor=blue,
    urlcolor=blue,
    }
\usepackage{array}
\usepackage{booktabs}
\usepackage{footnote}
\usepackage{longtable} % Nødvendig for at kunne inkludere tabeller
\makeatletter
\def\LT@err#1#2{}
\makeatother
\usepackage{caption}
\providecommand{\tightlist}{%
  \setlength{\itemsep}{0pt}%
  \setlength{\parskip}{0pt}%
}
\providecommand{\pandocbounded}[1]{#1}
$header-includes$
\usepackage{caption}
\usepackage{subcaption}
\usepackage{enumitem}
\usepackage{titlesec}
\setlist[itemize]{noitemsep, topsep=0pt, partopsep=0pt, parsep=1pt, left=0pt}
\setcounter{secnumdepth}{0}   % Som i Word-skabelonen: ingen nummerering. Sæt til 3 for at nummerere
\setcounter{tocdepth}{$if(toc-depth)$$toc-depth$$else$3$endif$}
\setlist[itemize,2]{label=\textopenbullet}
\setlist[enumerate,1]{leftmargin=36bp, labelindent=0pt, labelwidth=18bp,
  labelsep=0pt, align=left, itemsep=6bp, topsep=0pt, partopsep=0pt,
  parsep=0pt, first={\setlength{\rightskip}{3.5bp plus 1fil}}}
 
\usepackage{polyglossia}
\setdefaultlanguage{danish} % Dato på dansk
\usepackage{ragged2e}
 
\usepackage{amsmath}
\renewcommand{\thefootnote}{\arabic{footnote}}
\usepackage{fancyhdr}
 
% Sidetal nederst til højre (ikke på forsiden)
\fancypagestyle{rapport}{%
  \fancyhf{}
  \fancyfoot[R]{\thepage\hspace*{17.9bp}}
  \renewcommand{\headrulewidth}{0pt}
  \renewcommand{\footrulewidth}{0pt}
}
\pagestyle{rapport}
\setlength{\headheight}{14pt}
 
\newgeometry{top=105.5bp, bottom=3.2cm, left=2cm, right=2cm,
  headsep=0pt, footskip=25.1bp}
\setlength{\headwidth}{\textwidth}
\setlength{\parindent}{0pt}
\setlength{\parskip}{6bp}
 
% Som Word: ingen orphans/widows og ingen automatisk orddeling
\clubpenalty=10000
\widowpenalty=10000
 
% Brødtekst: Times New Roman 12 pt med 1,5 linjeafstand (Word: 20,7 pt)
\makeatletter
\renewcommand\normalsize{\@setfontsize\normalsize{12bp}{20.7bp}%
  \abovedisplayskip 12\p@ \@plus3\p@ \@minus7\p@
  \abovedisplayshortskip \z@ \@plus3\p@
  \belowdisplayshortskip 6.5\p@ \@plus3.5\p@ \@minus3\p@
  \belowdisplayskip \abovedisplayskip
  \let\@listi\@listI}
\makeatother
\normalsize
 
\definecolor{KUrod}{RGB}{144,26,30}
\definecolor{KUgra}{RGB}{102,102,102}
 
\setmainfont{Times New Roman}[Kerning=Off, Ligatures=NoCommon]
\setsansfont{Arial}[Kerning=Off, Ligatures=NoCommon]
 
% Overskrifter: Arial Bold 18 / 12 / 10 pt (afstande målt i Word-skabelonen)
\newif\ifkuafterh
\AddToHook{para/begin}{\global\kuafterhfalse}
\titleformat{\section}{\sffamily\bfseries\fontsize{18bp}{20.7bp}\selectfont}{\thesection\quad}{0pt}{}[\global\kuafterhtrue]
\titleformat{\subsection}{\sffamily\bfseries\fontsize{12bp}{13.8bp}\selectfont}{\thesubsection\quad}{0pt}{}
\titleformat{\subsubsection}{\sffamily\bfseries\fontsize{10bp}{11.5bp}\selectfont}{\thesubsubsection\quad}{0pt}{}
\titlespacing*{\section}{0pt}{17.9bp}{0.2bp}
\titlespacing*{\subsection}{0pt}{19bp}{-3bp}
\titlespacing*{\subsubsection}{0pt}{13.5bp}{-4.8bp}
% Word lægger afstandene sammen, når H2 følger direkte efter H1 (titlesec tager maksimum)
\let\kuoldsubsection\subsection
\renewcommand\subsection{\ifkuafterh\kern 9.9bp\global\kuafterhfalse\fi\kuoldsubsection}
 
% Indholdsfortegnelse: niveau 1 Arial Bold 12 versaler, niveau 2 TNR Bold 10, niveau 3 TNR 10
\makeatletter
\renewcommand*\numberline[1]{#1\hspace{0.8em}}
\newcommand*\kutocline[6]{% #1 skrift, #2 indrykning, #3 prik-afstand, #4 versaler (eller \@firstofone), #5 tekst, #6 side
  \par
  {\parindent=0pt \parskip=0pt \leftskip=#2 \rightskip=3em \parfillskip=-\rightskip
   \interlinepenalty\@M #1\leavevmode
   #4{#5}\nobreak
   \leaders\hbox to #3{\hss.}\hfill\nobreak
   \hbox to 2em{\hss #6}\par}}
\renewcommand*\l@section[2]{\kutocline
  {\sffamily\bfseries\fontsize{12bp}{37.2bp}\selectfont}{0bp}{3.36bp}\MakeUppercase{#1}{#2}}
\renewcommand*\l@subsection[2]{\kutocline
  {\rmfamily\bfseries\fontsize{10bp}{23.9bp}\selectfont}{0bp}{2.52bp}\@firstofone{#1}{#2}}
\renewcommand*\l@subsubsection[2]{\kutocline
  {\rmfamily\mdseries\fontsize{10bp}{17.3bp}\selectfont}{12bp}{2.52bp}\@firstofone{#1}{#2}}
\makeatother
 
\begin{document}
\thispagestyle{empty}
 
% ---------------- Forside ----------------
\begin{tikzpicture}[remember picture, overlay]
  % Rød streg i fuld sidebredde
  \draw[KUrod, line width=0.5bp]
    ([yshift=-211.15bp]current page.north west) --
    ([yshift=-211.15bp]current page.north east);
  % Segl (udsnit af KU_logo.png) og de to røde prikker
  \begin{scope}
    \clip ([xshift=507.55bp,yshift=-127.5bp]current page.north west)
      circle (33.5bp);
    \node[anchor=center, inner sep=0pt]
      at ([xshift=504.3bp,yshift=-156.6bp]current page.north west)
      {\includegraphics[width=133.4bp]{KU_logo.png}};
  \end{scope}
  \fill[KUrod] ([xshift=466.85bp,yshift=-176.4bp]current page.north west)
    circle (15.85bp);
  \fill[KUrod] ([xshift=507.55bp,yshift=-210.4bp]current page.north west)
    circle (8.4bp);
  % KØBENHAVNS UNIVERSITET
  \node[anchor=base west, inner sep=0pt]
    at ([xshift=55.4bp,yshift=-52.3bp]current page.north west)
    {\textls[196]{\fontsize{11bp}{11bp}\selectfont KØBENHAVNS UNIVERSITET}};
  % Titel og undertitel
  \node[anchor=base west, inner sep=0pt]
    at ([xshift=63.9bp,yshift=-274.8bp]current page.north west)
    {\parbox[t]{460bp}{\raggedright\hyphenpenalty=10000
      \sffamily\bfseries\fontsize{24bp}{27.6bp}\selectfont \Titel\par
      \vspace{36bp}%
      \mdseries\fontsize{18bp}{20.7bp}\selectfont \Undertitel\par
      $if(forsidebillede)$\vspace{10bp}%
      \includegraphics[width=\linewidth]{$forsidebillede$}$endif$}};    
  % Grafik: to tynde grå streger, der løber ud over sidekanten
  \begin{scope}
    \clip (current page.north west) rectangle (current page.south east);
    \draw[KUrod, line width=0.6bp]
      ([xshift=620bp,yshift=-300bp]current page.north west) --
      ([xshift=380bp,yshift=-930bp]current page.north west);
    \draw[KUrod, line width=0.6bp]
      ([xshift=640bp,yshift=-520bp]current page.north west) --
      ([xshift=330bp,yshift=-930bp]current page.north west);
  \end{scope}
  % Afsender (over datoen)
  \node[anchor=base west, inner sep=0pt]
    at ([xshift=63.9bp,yshift=-700.2bp]current page.north west)
    {\sffamily\fontsize{14bp}{14bp}\selectfont \Afsender};
  % Dato
  \node[anchor=base west, inner sep=0pt]
    at ([xshift=63.9bp,yshift=-721.2bp]current page.north west)
    {\sffamily\fontsize{14bp}{14bp}\selectfont \Dato};
\end{tikzpicture}
\clearpage
 
% ---------------- Indholdsfortegnelse ----------------
$if(toc)$
\begin{tikzpicture}[remember picture, overlay]
  \node[anchor=base west, inner sep=0pt]
    at ([xshift=56.7bp,yshift=-107.5bp]current page.north west)
    {\sffamily\bfseries\fontsize{24bp}{27.6bp}\selectfont \Indholdstitel};
\end{tikzpicture}%
\vspace*{-12.4bp}
{\hypersetup{linkcolor=black}\makeatletter\@starttoc{toc}\makeatother}
\clearpage
$endif$
 
% ---------------- Brødtekst ----------------
\begingroup
\hyphenpenalty=10000
\rightskip=3.5bp plus 1fil\relax
\microtypesetup{protrusion=false}
$body$
\par
\endgroup
 
\end{document}
